\begin{lemma}
Under the hypotheses of the opposite-projection collision bound, the probability that the size bound holds but the blue estimate fails is $o(1)$. Hence the conditional blue estimate holds with high probability.
\end{lemma}
\begin{proof}
Write \(M=\noderef{TwoBiteNaturalReducedVertexCount}(n)\),
\(p=\noderef{TwoBiteEdgeProbability}(1/2,n)\), and \(L=\log n\).  Fix a
blue ordered pair \(b\ne c\).  Let \(E_{b,c}\) be the event that the size
bound holds and this ordered pair violates the claimed blue projected
intersection estimate.

We average over blue conditioning cells using
\(\noderef{OppositeProjectionBlueConditioning}\).  For all sufficiently large
\(n\), \(0\le p\le1\): indeed \(L\ge0\) eventually, so \(p\ge0\), and
\(\sqrt{L/n}\to0\), so \(p\le1\) eventually.  Therefore
\(\noderef{TwoBiteSampleWeightNonnegative}\) gives nonnegative sample
weights.  It is enough to prove, uniformly for every blue base graph \(G_B\)
and every blue row map \(\rho:\operatorname{Fin}n\to\operatorname{Fin}M\),
that
\[
\operatorname{TwoBiteConditionalProbability}
  (E_{b,c}\mid \omega.2.1=G_B,\ \pi_B=\rho)
\le \exp(-cL^3)
\]
for a constant \(c>0\).  The helper then multiplies these conditional bounds
by the masses of the finite blue cells, sums over the cells, handles
zero-mass cells by the conditional-probability zero convention, and uses the
total mass bound to obtain
\[
\operatorname{TwoBiteEventProbability}(E_{b,c})\le \exp(-cL^3).
\]

Fix such a blue cell with nonzero mass.  In this cell the blue base graph and
blue row map are fixed, so the lifted blue neighborhoods of \(b\) and \(c\)
are fixed subsets of original vertices.  Put
\[
U=N_b,\qquad V=N_c,\qquad U_0=U\setminus V,\qquad V_0=V\setminus U.
\]
Zero-mass cells need not be considered further, since they contribute
nothing to the average in \(\noderef{OppositeProjectionBlueConditioning}\).

On the size event,
\[
|U_0|\le |U|\le 2pn,\qquad |V_0|\le |V|\le 2pn.
\]
The fixed-exposure tail estimate needed for the remaining red coordinates is
packaged in \(\noderef{OppositeProjectionBlueExposureTailBound}\).  That
helper expands the eta-conditioned finite sample weights, handles zero-mass
conditioning cells, and proves directly by rowwise injection counting that if
\(T\subseteq V_0\) is fixed and \(S\) is a fixed set of red coordinates, then
the conditional probability that every vertex of \(T\) lands in \(S\) is at
most \((|S|/M)^{|T|}\).

Now expose all red coordinates of \(U_0\).  Let \(\eta\) denote a possible
full exposure of these coordinates, and write
\[
S(\eta)=\pi_R(U_0),\qquad |S(\eta)|\le |U_0|.
\]
The rows used by \(U_0\) and \(V_0\) are disjoint: \(U_0\) lies over blue
rows adjacent to \(b\) but not \(c\), while \(V_0\) lies over blue
rows adjacent to \(c\) but not \(b\).  The sets are also disjoint by definition of
the set differences.  Thus the hypotheses of
\(\noderef{OppositeProjectionBlueExposureTailBound}\) hold for every exposure
cell \(\eta\).

For a fixed exposure \(\eta\), define
\[
W_\eta=|\{v\in V_0:\pi_R(v)\in S(\eta)\}|.
\]
The exposure-tail helper first handles the zero-mass exposure cells using the
zero convention in \(\noderef{TwoBiteConditionalProbability}\).  In every
nonzero cell it expands the finite sample weights, cancels the row factors
fixed by the \(U_0\)-exposure, and performs the color-swapped rowwise
injection count on the rows meeting \(V_0\).  Consequently, for every \(t\),
\[
\Pr(W_\eta\ge t\mid G_B,\rho,\eta)
\le \binom{|V_0|}{t}\left(\frac{|S(\eta)|}{M}\right)^t
\le \binom{|V_0|}{t}\left(\frac{|U_0|}{M}\right)^t .
\]
Average this inequality over all possible exposures \(\eta\).  The exposure
events form a finite partition of the already conditioned cell; each has
nonnegative mass, and zero-mass exposure events contribute \(0\).  Since the
right hand side is uniform in \(\eta\), the same bound holds before exposure:
\[
\Pr(W\ge t\mid G_B,\rho)
\le \binom{|V_0|}{t}\left(\frac{|U_0|}{M}\right)^t
\le \left(\frac{e|U_0||V_0|}{Mt}\right)^t.
\]
The last inequality is the standard estimate
\(\binom bt\le(eb/t)^t\), and the event is empty when \(t>|V_0|\).

The intersection-containment node
\(\noderef{OppositeProjectionIntersectionContainment}\) gives
\[
|\pi_R(U)\cap\pi_R(V)|
\le |U\cap V|+W.
\]
Indeed, the only projected intersection points not already coming from
\(U\cap V\) are witnessed by a vertex of \(V_0\) whose red coordinate lies in
\(\pi_R(U_0)\), and these witnesses are counted by \(W\).

Now estimate the parameter in the tail bound.  Since
\[
p=\noderef{TwoBiteEdgeProbability}(1/2,n)
=\frac12\sqrt{\frac{L}{n}},
\]
and \(L\ge0\) for all sufficiently large \(n\), we have
\[
p^2n=\frac14\frac{L}{n}\,n=\frac{L}{4}.
\]
Also
\[
M=\left\lfloor\frac{n}{L^2}\right\rfloor .
\]
Because \(n/L^2\to\infty\), for all sufficiently large \(n\) we have
\(n/L^2\ge2\), and then
\[
M=\left\lfloor\frac{n}{L^2}\right\rfloor
\ge \frac12\frac{n}{L^2}.
\]
Indeed, for \(x\ge2\), \(\lfloor x\rfloor\ge x-1\ge x/2\).  Combining this
with the size bounds gives, for all sufficiently large \(n\),
\[
\frac{|U_0||V_0|}{M}
\le \frac{(2pn)^2}{M}
=\frac{4p^2n^2}{M}
\le \frac{Ln}{M}
\le 2L^3.
\]
Take \(t=\lceil100L^3\rceil\).  Then \(t\ge100L^3\), so
\[
\frac{e|U_0||V_0|}{Mt}
\le \frac{2eL^3}{100L^3}
=\frac e{50}.
\]
Since \(e/50<1\), the tail estimate gives
\[
\left(\frac{e|U_0||V_0|}{Mt}\right)^t
\le \left(\frac e{50}\right)^t
\le \left(\frac e{50}\right)^{100L^3}
=\exp(-cL^3),
\]
where \(c=100\log(50/e)>0\).  If \(t>|V_0|\), the tail event is empty, so the
same upper bound holds trivially.  Thus a fixed blue pair fails with
probability at most \(\exp(-cL^3)\), uniformly over the conditioning data.

There are at most \(M^2\le n^2=\exp(2L)\) ordered blue pairs for all large
\(n\).  The union bound gives
\[
M^2\exp(-cL^3)\to0
\]
by \(\noderef{OppositeProjectionAsymptoticBound}\).  Therefore, with high
probability, no blue pair violates the desired projected intersection estimate
while the size bound holds.
\end{proof}
